\chapter{Stratégies de futures à court terme}\label{ch:s1:short-term-futures-strategies}

Lucca et Moench ont relevé de forts rendements excédentaires moyens sur les actions américaines dans les heures qui précèdent les décisions de politique monétaire de la Réserve fédérale~: des rendements
qui ont augmenté au fil du temps, représentent une part notable du rendement annuel du marché, apparaissent dans les indices d'autres pays et n'ont aucun équivalent avant les autres
annonces macroéconomiques. Une fois publié, un tel effet appartient à tout le monde. Les stratégies sur futures à court terme vivent de régularités de ce type, au sein de la
journée de négociation~: la première demi-heure, les mouvements qui rebondissent en quelques minutes, les heures qui précèdent une publication programmée. Sur les séances synthétiques de ce chapitre
la détention avant annonce rapporte 40.6 points de base par événement avant l'encombrement injecté et 18.8 après, avec un ratio de Sharpe qui passe de 1.21 à
0.58~; une contre-tendance sur les grands mouvements de cinq minutes rapporte un ratio de Sharpe de 2.86 avant coûts et $-1.50$ après une demi-fourchette dans chaque sens. La construction est
\texttt{firm.futintraday}.

\section{Sessions en barres}

La séance régulière du contrat à terme sur indice est simulée en barres d'une minute (le \cref{lst:s1:short-term-futures-strategies:sessions})~: dix ans de séances
de 390 minutes, une volatilité annuelle de 16.2~\%, avec un quart de la variance dans le gap overnight, un profil intrajournalier en U (deux fois plus volatil à l'ouverture et
à la clôture qu'à midi), une dérive quotidienne de tendance d'écart-type 0.4~\% sur la séance, et un rebond de $-0.1$ dans les rendements à une minute. Huit fois par
an, une annonce programmée survient à la minute 270~; une dérive de 0.5~\% est injectée entre le midi de la veille et l'annonce, et à partir de la sixième
année elle est réduite à 30~\% de ce niveau~: l'encombrement qui suit la publication. Le \texttt{firm.tape} du livre~7, au niveau des messages, et le \texttt{firm.evbt}, piloté par les événements, sont
les outils pour les questions au niveau de la file~; pour dix ans de séances, les barres suffisent, et le tape fournit le coût. Sur une séance simulée du tape (un cours de
100.00, un tick de 0.01, donc un tick vaut un point de base) la fourchette cotée pondérée par le temps vaut 1.08 ticks~; chaque transaction paie ici la moitié de celle-ci, soit 0.54 point de base.

\section{Plages d'ouverture et retour à la moyenne intrajournalier}

\begin{definition}[Cassure de la plage d'ouverture]\label{def:s1:short-term-futures-strategies:orb}
Une \emph{cassure de la plage d'ouverture}\index{cassure de la plage d'ouverture} enregistre les plus hauts et les plus bas des premières minutes de la séance (la
plage d'ouverture) et prend une position dans le sens du premier mouvement ultérieur qui la franchit, tenue jusqu'à la clôture ou jusqu'à un stop.
\end{definition}

\begin{definition}[Retour à la moyenne intrajournalier]\label{def:s1:short-term-futures-strategies:mr}
Le \emph{retour à la moyenne intrajournalier}\index{retour à la moyenne intrajournalier} est la tendance des mouvements de cours intrajournaliers courts à être en partie annulés en quelques minutes, sous l'effet
du rebond bid-ask et de l'inventaire des fournisseurs de liquidité~; une stratégie de retour à la moyenne prend le contre-pied des mouvements supérieurs à un seuil et sort au bout d'une durée de détention fixe.
\end{definition}

Les deux règles misent sur des caractéristiques opposées de la séance synthétique (le \cref{lst:s1:short-term-futures-strategies:rules}). La cassure parie sur la journée de tendance~:
dès que la séance sort de la plage de ses premières minutes, la dérive de la journée a plus de chances d'aller dans ce sens. La contre-tendance parie sur le rebond~: une fois que
le cours s'est écarté de plus de $z$ écarts-types en cinq minutes, on prend le contre-pied du mouvement et l'on sort cinq minutes plus tard.

\begin{center}
\small
\begin{tabular}{@{}p{7.0cm}ccc@{}}
\toprule
plage d'ouverture (minutes) & 15 & 30 & 60\\
\midrule
cassure~: ratio de Sharpe avant coûts & 0.56 & 0.77 & 1.49\\
cassure~: ratio de Sharpe après coûts & 0.35 & 0.55 & 1.25\\
cassure~: rendement quotidien moyen après coûts (pb) & 1.88 & 2.77 & 5.62\\
\bottomrule
\end{tabular}
\end{center}

Une plage plus longue donne une lecture plus tardive mais plus sûre de la direction de la journée~: la dérive injectée se répartit uniformément sur la séance, si bien qu'une heure en révèle davantage
qu'un quart d'heure, et la cassure ne prend position qu'une fois par jour au plus, de sorte que les coûts pèsent peu. La contre-tendance est l'inverse (la \cref{fig:s1:short-term-futures-strategies:fade})~:
le rebond est assez fort pour donner des ratios de Sharpe de 0.59 à 3.78 avant coûts, mais à 0.54 point de base dans chaque sens, chaque seuil est perdant. Le rebond rend
un dixième du mouvement de la dernière minute, une petite part d'un mouvement de cinq minutes, et cette part est plus petite qu'une fourchette.

\begin{omfigure}
\begin{tikzpicture}
\begin{axis}[width=10cm,height=5cm,ybar,bar width=12pt,xlabel={\small seuil (écarts-types d'un mouvement de cinq minutes)},
  ylabel={\small ratio de Sharpe},tick label style={font=\footnotesize},symbolic x coords={1.5,2.0,2.5,3.0},xtick=data,ymin=-4,ymax=4.5,
  enlarge x limits=0.2,grid=major,grid style={gray!25},legend columns=2,legend style={font=\scriptsize,at={(0.5,-0.3)},anchor=north,
  draw=none,fill=none,/tikz/every even column/.append style={column sep=8pt}},area legend,table/col sep=comma]
\addplot[fill=omExo!50,draw=omExo] table[x=z,y=gross]{figdata/strategies-1/25-short-term-futures-strategies/fade.csv};
\addplot[fill=omThm!60,draw=omThm] table[x=z,y=net]{figdata/strategies-1/25-short-term-futures-strategies/fade.csv};
\legend{avant coûts,après une demi-fourchette dans chaque sens}
\end{axis}
\end{tikzpicture}

\omcaption{La contre-tendance intrajournalière sur dix ans de séances synthétiques à une minute~: ratios de Sharpe avant et après coûts d'une contre-tendance sur les mouvements de cinq minutes au-delà
d'un seuil et en la conservant cinq minutes. Le rebond injecté est réel~; la demi-fourchette cotée du tape du livre~7, dans chaque sens, est plus grande. Données~:
\texttt{s1\_futures.fade\_table}.}\label{fig:s1:short-term-futures-strategies:fade}
\end{omfigure}

\section{Annonces programmées}

\begin{definition}[Dérive pré-annonce]\label{def:s1:short-term-futures-strategies:pre}
Une \emph{dérive pré-annonce}\index{dérive pré-annonce} est un mouvement moyen des cours dans les heures qui précèdent une annonce programmée, dans une direction
qui ne dépend pas du contenu de l'annonce~; la dérive pré-FOMC sur les actions américaines est le cas le plus connu.
\end{definition}

\begin{dated}{2026-09}{Réunions du FOMC en 2026}\label{dat:s1:short-term-futures-strategies:fomc}
Le calendrier de la Réserve fédérale recense huit réunions du FOMC en 2026~: 27--28 janvier, 17--18 mars, 28--29 avril, 16--17 juin, 28--29 juillet, 15--16
septembre, 27--28 octobre et 8--9 décembre~; les réunions de mars, juin, septembre et décembre comprennent un Summary of Economic Projections.
\end{dated}

La dérive de Lucca et Moench est singulière à deux égards~: elle précède l'information au lieu de la suivre, et elle a un signe quelle que soit l'annonce. Elle
ne se manifeste aussi que sur les actions~; ils n'en ont trouvé ni sur les Treasuries ni sur les contrats à terme de marché monétaire, ni avant les autres annonces majeures. Cela en fait une
prime de risque ou une régularité comportementale, et non une information, et elle est facile à exploiter~: il suffit de détenir le contrat à terme sur indice actions depuis la veille jusqu'à
l'annonce.

\begin{center}
\small
\begin{tabular}{@{}lcc@{}}
\toprule
détention avant annonce, après coûts & années 1--5 & années 6--10 (encombrement)\\
\midrule
événements & 39 & 40\\
rendement moyen par événement (pb) & 40.6 & 18.8\\
statistique t & 2.7 & 1.3\\
ratio de Sharpe annuel (huit événements) & 1.21 & 0.58\\
\bottomrule
\end{tabular}
\end{center}

La dérive injectée vaut 50 points de base avant l'encombrement et 15 après. Les moyennes mesurées, 40.6 et 18.8, sont à moins d'une erreur type de ces valeurs~: avec une
volatilité journalière d'environ 1~\% par événement, 39 événements mesurent la moyenne à environ 16 points de base (la \cref{fig:s1:short-term-futures-strategies:path}). C'est
la difficulté pratique des stratégies événementielles~: huit événements par an forment un petit échantillon, et une décroissance au tiers de l'effet n'est guère visible en cinq
années. La dérive publiée constitue un a priori fort, et un portefeuille réel ne peut la confirmer que lentement.

\begin{omfigure}
\begin{tikzpicture}
\begin{axis}[width=11cm,height=5.2cm,xlabel={\small minutes depuis le midi de la veille},ylabel={\small rendement cumulé moyen (pb)},
  tick label style={font=\footnotesize},xmin=0,xmax=590,ymin=-20,ymax=70,grid=major,grid style={gray!25},legend pos=north west,
  legend style={font=\scriptsize},table/col sep=comma]
\addplot[omThm,thick] table[x=minute,y=before]{figdata/strategies-1/25-short-term-futures-strategies/announcement.csv};
\addplot[omDef,thick] table[x=minute,y=after]{figdata/strategies-1/25-short-term-futures-strategies/announcement.csv};
\addplot[black!50,dashed,forget plot] coordinates {(466,-20) (466,70)};
\legend{années 1--5,années 6--10 (encombrement)}
\end{axis}
\end{tikzpicture}

\omcaption{Trajectoire moyenne du contrat à terme sur indice synthétique autour des annonces programmées, depuis le midi de la veille (l'écart overnight compté comme
une étape), jusqu'à la clôture du jour de l'annonce~; la ligne en tirets marque l'annonce. Effet injecté~: 50 points de base avant l'annonce pendant les cinq premières années, 15 pendant
les cinq dernières. Données~: \texttt{s1\_futures.announcement\_path}.}\label{fig:s1:short-term-futures-strategies:path}
\end{omfigure}

\section{Ce qui fait survivre un edge à court terme}

Trois caractéristiques tranchent. Rotation contre edge~: le fade traite environ sept fois par jour, pour une fraction de point de base à chaque fois, et perd~; la cassure traite
une seule fois et conserve l'essentiel de son edge. Taille d'échantillon face à l'érosion~: une stratégie événementielle à huit événements par an ne peut pas voir son propre encombrement pendant des années. Et le
mécanisme~: un rebond est réel, mais il revient aux fournisseurs de liquidité qui touchent la fourchette~; la direction d'une journée de tendance et une dérive programmée sont plus difficiles
à expliquer, et donc plus délicates à confirmer. L'écart overnight et la réaction après les annonces sont d'autres pistes~; les fiches de stratégie les consignent
sans chiffre de performance vérifié.

\section{Fiches de stratégie}

\begin{strategyfile}{Cassure de la plage d'ouverture}\label{strat:s1:short-term-futures-strategies:orb}
\sfield{Qui paie, et pourquoi} Des traders qui réagissent tard à la direction de la journée.
\sfield{Instruments et plateformes} Contrats à terme liquides sur indices, taux et matières premières.
\sfield{Signal} Le premier mouvement en dehors de la plage d'ouverture.
\sfield{Taille et exécution} Une position par jour, conservée jusqu'à la clôture ou un stop~; dimensionnée selon la volatilité.
\sfield{Coûts} Faibles~: un aller-retour par jour.
\sfield{Comment elle s'éteint} Journées sans tendance~; encombrement aux bornes de la plage.
\sfield{Horizon, capacité, infrastructure} Quelques heures~; barres intrajournalières.
\sfield{Backtest en toute honnêteté} Exécutions au niveau de la cassure, auxquelles s'ajoute un glissement~; longueur de la plage choisie avant le test.
\sfield{Sources} Aucun chiffre de performance vérifié~; ce chapitre~: 0.55 après coûts avec une plage de 30 minutes, synthétique.
\end{strategyfile}

\begin{strategyfile}{Retour à la moyenne intrajournalier sur contrats à terme}\label{strat:s1:short-term-futures-strategies:mr}
\sfield{Qui paie, et pourquoi} Des traders impatients qui poussent brièvement les cours~; en pratique, ce sont les fournisseurs de liquidité qui le perçoivent.
\sfield{Instruments et plateformes} Contrats à terme liquides.
\sfield{Signal} Mouvements sur cinq minutes au-delà d'un seuil.
\sfield{Taille et exécution} Fade, tenue de quelques minutes~; entrée passive si possible.
\sfield{Coûts} Décisifs~: les ordres agressifs paient la fourchette sur chaque transaction.
\sfield{Comment elle s'éteint} Coûts~; fournisseurs de liquidité plus rapides.
\sfield{Horizon, capacité, infrastructure} Minutes~; exécution tenant compte de la position dans la file (le backtester de niveau 2 du livre~7).
\sfield{Backtest en toute honnêteté} La fourchette sur chaque exécution agressive~; la position dans la file pour les exécutions passives.
\sfield{Sources} Ce chapitre~: 2.86 avant coûts, $-1.50$ après, synthétique.
\end{strategyfile}

\begin{strategyfile}{Dérive pré-FOMC}\label{strat:s1:short-term-futures-strategies:prefomc}
\sfield{Qui paie, et pourquoi} Incertain~: une prime pour porter le risque d'annonce, ou le comportement des investisseurs avant les nouvelles de politique monétaire.
\sfield{Instruments et plateformes} Contrats à terme sur indices actions.
\sfield{Signal} Le calendrier du FOMC.
\sfield{Taille et exécution} Position longue, de la veille jusqu'à l'annonce.
\sfield{Coûts} Un aller-retour par réunion.
\sfield{Comment elle s'éteint} Publication et encombrement.
\sfield{Horizon, capacité, infrastructure} Une journée, huit fois par an~; grande capacité.
\sfield{Backtest en toute honnêteté} Dates de réunion et heures des communiqués telles que programmées~; réunions non programmées exclues~; événements suffisants pour juger de l'érosion.
\sfield{Sources} Lucca et Moench (2015)~; ce chapitre~: 40.6 points de base par événement avant encombrement, 18.8 après.
\end{strategyfile}

\begin{strategyfile}{Momentum post-annonce}\label{strat:s1:short-term-futures-strategies:post}
\sfield{Qui paie, et pourquoi} Des traders qui s'ajustent progressivement aux nouvelles programmées.
\sfield{Instruments et plateformes} Contrats à terme sur taux et indices autour des publications de données.
\sfield{Signal} La direction des premières minutes après une publication, par rapport à la surprise.
\sfield{Taille et exécution} Suivre le mouvement, tenue de plusieurs heures.
\sfield{Coûts} Fourchettes élargies lors des publications.
\sfield{Comment elle s'éteint} Des algorithmes qui achèvent la réaction en quelques secondes (chapitre~17).
\sfield{Horizon, capacité, infrastructure} De quelques minutes à quelques heures~; données de publication à faible latence.
\sfield{Backtest en toute honnêteté} Horodatages des publications à la seconde près~; fourchettes au moment de la publication.
\sfield{Sources} Aucun chiffre de performance vérifié.
\end{strategyfile}

\begin{strategyfile}{Fade de l'écart overnight}\label{strat:s1:short-term-futures-strategies:gap}
\sfield{Qui paie, et pourquoi} Des ordres passés overnight qui s'emballent à l'ouverture.
\sfield{Instruments et plateformes} Contrats à terme sur indices à l'ouverture du marché au comptant.
\sfield{Signal} L'écart entre la clôture précédente et l'ouverture, rapporté à sa volatilité.
\sfield{Taille et exécution} Fade des écarts importants, sortie dans la matinée.
\sfield{Coûts} Fourchettes d'ouverture élargies.
\sfield{Comment elle s'éteint} Les écarts qui portent une nouvelle information se prolongent.
\sfield{Horizon, capacité, infrastructure} Quelques heures.
\sfield{Backtest en toute honnêteté} Le cours d'ouverture effectivement négociable, et non le dernier cours échangé~; jours d'actualité séparés.
\sfield{Sources} Aucun chiffre de performance vérifié~; le \omterm{def:s1:intraday-patterns:returns}{rendement overnight} du chapitre~13.
\end{strategyfile}

\section{Tutoriel~: avant l'annonce}\label{tut:s1:short-term-futures-strategies:tut}

\begin{tutorial}
\textbf{Objectif.} Simuler dix ans de sessions de contrats à terme à la minute, avec une dérive de journée de tendance, un rebond et une \omterm{def:s1:short-term-futures-strategies:pre}{dérive pré-annonce} encombrée~; prendre le coût
d'une transaction à partir du tape du livre~7~; tester la cassure, le fade et la tenue pré-annonce. \textbf{État final~:} les tableaux et les deux figures.
\begin{tutsteps}
\item \textbf{Sessions}~: le profil en U, le rebond, la journée de tendance et les annonces injectées.
\omcode{code/firm/futintraday/firm_futintraday.py}{50}{71}{Dix ans de sessions à la minute.}\label{lst:s1:short-term-futures-strategies:sessions}
\item \textbf{Les règles}~: la \omterm{def:s1:short-term-futures-strategies:orb}{cassure de la plage d'ouverture} et le fade.
\omcode{code/firm/futintraday/firm_futintraday.py}{74}{107}{Cassure de la plage d'ouverture et fade intrajournalier.}\label{lst:s1:short-term-futures-strategies:rules}
\item \textbf{Exécuter} \texttt{tape\_cost()}, \texttt{orb\_table()}, \texttt{fade\_table()}, \texttt{announcements()} et \texttt{announcement\_path()}, et
\texttt{fig\_futures.py}.
\end{tutsteps}
\textbf{Pour aller plus loin.} Passer le fade de façon passive, du côté éloigné de la fourchette, et modéliser la probabilité d'exécution~; faire tourner la stratégie pré-annonce sur
quarante ans d'événements pour voir combien de temps il faut pour détecter l'encombrement~; remplacer les barres par des sessions firm.tape pour les minutes d'annonce.
\end{tutorial}

\section{Construction~: boîte à outils des contrats à terme intrajournaliers}\label{bld:s1:short-term-futures-strategies:futintraday}

\begin{build}
\bfield{Rôle} Sessions de contrats à terme au niveau des barres, avec effets injectés, et les stratégies de cassure, de fade et de pré-annonce avec coûts.
\bfield{Interface} \texttt{SessionConfig(\dots)}, \texttt{simulate\_sessions(cfg, rng)}, \texttt{orb(S, window, cost)},
\texttt{fade(S, lookback, z, hold, cost)}, \texttt{pre\_announcement(S, cost)}.
\bfield{Règles} Décisions prises à partir de barres déjà clôturées~; une position à la fois~; deux coûts par aller-retour.
\bfield{Tests d'acceptation} \texttt{code/firm/futintraday/tests/}~: la volatilité et le rebond des sessions, ainsi qu'une cassure et une tenue pré-annonce vérifiées à la main.
\bfield{Pour aller plus loin} Exécutions passives~; stratégies sur écarts~; annonces avec surprises et dynamiques post-annonce.
\end{build}

\begin{omsources}
\item D.~O.~Lucca et E.~Moench, ``The pre-FOMC announcement drift'', \emph{Journal of Finance} 70(1), 2015.
\item Board of Governors of the Federal Reserve System, calendriers des réunions du FOMC.
\end{omsources}

\section{Exercices}

\begin{exercise}[$\star$]\label{exo:s1:short-term-futures-strategies:1}
Au prix de 100.00 et avec un tick de 0.01, combien de points de base vaut un tick~? Quel est le coût d'un aller-retour pour une fourchette cotée de 1.08 tick~?
\end{exercise}

\begin{exercise}[$\star$]\label{exo:s1:short-term-futures-strategies:2}
La dérive plantée de 50 points de base est ramenée à 30~\%. Que reste-t-il, et quel multiplicateur annuel le fait de trader huit événements par an
donne-t-il au ratio de Sharpe par événement~?
\end{exercise}

\begin{exercise}[$\star$]\label{exo:s1:short-term-futures-strategies:3}
Pourquoi une plage d'ouverture plus longue fonctionne-t-elle mieux ici~?
\end{exercise}

\begin{exercise}[$\star\star$]\label{exo:s1:short-term-futures-strategies:4}
Avec un écart-type par événement d'environ 1~\%, quelle est l'erreur type d'une moyenne sur 39 événements, et quelle statistique t une vraie dérive de 50
points de base produit-elle en moyenne~?
\end{exercise}

\begin{exercise}[$\star\star$]\label{exo:s1:short-term-futures-strategies:5}
Pourquoi le fade perd-il après coûts à chaque seuil, alors qu'il gagne avant coûts~?
\end{exercise}

\begin{exercise}[$\star\star$]\label{exo:s1:short-term-futures-strategies:6}
Un rebond de $-0.1$ dans les rendements à une minute réduit la variance de session d'une journée. De combien réduit-il la volatilité d'une journée lorsque la session porte les trois
quarts de la variance~?
\end{exercise}

\begin{exercise}[$\star\star\star$]\label{exo:s1:short-term-futures-strategies:7}
\emph{Programmation.} Exécuter \texttt{fade\_stats(2.0)}. Combien le fade rapporte-t-il par jour avant coûts, combien de transactions réalise-t-il, et quel coût par côté
l'amènerait-il au seuil de rentabilité~?
\end{exercise}

\begin{exercise}[$\star\star\star$]\label{exo:s1:short-term-futures-strategies:8}
\emph{Trouver la faille.} «~La dérive pré-FOMC a rapporté 40 points de base par réunion au cours des cinq dernières années, donc l'avantage est intact.~»
\end{exercise}

\section{Problème~: avant l'annonce}

\begin{problem}[{Problème du week-end --- trois avantages intrajournaliers}]\label{pb:s1:short-term-futures-strategies:1}
Les sessions synthétiques du chapitre et le registre public.

\medskip\noindent\textbf{Partie I --- Sessions.}
\begin{enumerate}
\item Décrire le modèle de barres et ce qui y est planté.
\item Comment obtient-on le coût d'une transaction à partir de la bande du livre~7~?
\item Définir la \omterm{def:s1:short-term-futures-strategies:orb}{cassure de la plage d'ouverture} et le \omterm{def:s1:short-term-futures-strategies:mr}{retour à la moyenne intrajournalier}.
\item À quelle caractéristique de la session chacun de ces paris s'adosse-t-il~?
\end{enumerate}

\medskip\noindent\textbf{Partie II --- Règles.}
\begin{enumerate}[resume]
\item Donner les ratios de Sharpe de la cassure selon la longueur de la plage.
\item Donner les ratios de Sharpe du fade avant et après coûts.
\item Pourquoi l'un survit-il aux coûts et non l'autre~?
\item Qui encaisse le rebond sur les marchés réels~?
\end{enumerate}

\medskip\noindent\textbf{Partie III --- Annonces.}
\begin{enumerate}[resume]
\item Définir la \omterm{def:s1:short-term-futures-strategies:pre}{dérive pré-annonce}~; que mettent en évidence Lucca et Moench~?
\item Que liste l'encadré daté~?
\item Donner les rendements par événement et les ratios de Sharpe avant et après encombrement (crowding).
\item Pourquoi l'encombrement est-il difficile à détecter~?
\end{enumerate}

\medskip\noindent\textbf{Partie IV --- Le verdict.}
\begin{enumerate}[resume]
\item Énoncer le \emph{résultat nommé}~: le ratio de Sharpe de la \omterm{def:s1:short-term-futures-strategies:pre}{dérive pré-annonce} avant et après l'encombrement planté.
\item Énumérer les trois caractéristiques qui décident si un avantage de court terme survit.
\item Comment backtester honnêtement la dérive pré-FOMC~?
\item Quelle fiche de stratégie dépend le plus de la latence~?
\item Quel lien le chapitre entretient-il avec les motifs intrajournaliers du chapitre~13~?
\item De quoi faut-il disposer pour exécuter le fade avec profit~?
\item Pourquoi la dérive pré-FOMC n'est-elle pas une information~?
\item En une phrase~: que vendent les stratégies sur contrats à terme de court terme~?
\end{enumerate}
\end{problem}

\section{Questions d'entretien}

\begin{interviewq}[$\star$ \iqroles{trader}]\label{iq:s1:short-term-futures-strategies:1}
Qu'est-ce qu'une \omterm{def:s1:short-term-futures-strategies:orb}{cassure de la plage d'ouverture}, et quand échoue-t-elle~?
\end{interviewq}

\begin{interviewq}[$\star\star$ \iqroles{chercheur}]\label{iq:s1:short-term-futures-strategies:2}
Comment testeriez-vous si une \omterm{def:s1:short-term-futures-strategies:pre}{dérive pré-annonce} publiée existe encore~?
\end{interviewq}

\begin{interviewq}[$\star\star$ \iqroles{chercheur}]\label{iq:s1:short-term-futures-strategies:3}
Votre stratégie intrajournalière affiche un ratio de Sharpe de 3 avant coûts. Que vérifiez-vous en premier~?
\end{interviewq}

\begin{interviewq}[$\star\star$ \iqroles{développeur}]\label{iq:s1:short-term-futures-strategies:4}
De quoi un backtester a-t-il besoin pour simuler fidèlement une stratégie intrajournalière sur contrats à terme~?
\end{interviewq}

\begin{interviewq}[$\star\star$ \iqroles{risque}]\label{iq:s1:short-term-futures-strategies:5}
Quels sont les risques de détenir des contrats à terme sur indice à l'approche d'annonces programmées~?
\end{interviewq}

\begin{interviewq}[$\star\star\star$ \iqroles{chercheur}]\label{iq:s1:short-term-futures-strategies:6}
Les rendements à une minute s'écrivent $r_t = e_t + \theta e_{t-1}$, avec des $e_t$ indépendants de variance $\sigma^2$. Un trader parie contre le rendement de la dernière minute et conserve la
position une minute. Calculer le P\&L espéré par transaction et le coût par transaction au seuil de rentabilité.
\end{interviewq}
